\documentclass[11pt, a4paper]{article}

% Paquetes de idioma y tipografía
\usepackage[utf8]{inputenc}
\usepackage[spanish, es-tabla]{babel}
\usepackage[T1]{fontenc}
\usepackage{mathpazo} % Tipografía Palatino
\usepackage{microtype}

% Geometría de la página
\usepackage[margin=2.5cm, headheight=15pt]{geometry}

% Paquetes matemáticos
\usepackage{amsmath, amssymb, amsfonts, bm}

% Gráficos, colores y tablas
\usepackage{graphicx}
\usepackage[table, xcdraw]{xcolor}
\usepackage{booktabs}
\usepackage{tabularx}
\usepackage[most]{tcolorbox}
\usepackage{tikz}
\usetikzlibrary{shapes.geometric, arrows.meta, positioning, calc}

% Personalización de listas
\usepackage{enumitem}

% Encabezados y pies de página
\usepackage{fancyhdr}
\pagestyle{fancy}
\fancyhf{}
\lhead{\textbf{UCB Sede Tarija} | Ing. Mecatrónica}
\chead{Robótica (IMT-342)}
\rhead{Clase 1 - Sem 2 (V2)}
\lfoot{Prof: M.Sc. Ing. Bernardo Quiroga Turdera}
\rfoot{Página \thepage}
\renewcommand{\headrulewidth}{0.4pt}
\renewcommand{\footrulewidth}{0.4pt}

% Definición de colores institucionales
\definecolor{ucbblue}{RGB}{0, 51, 102}
\definecolor{ucbyellow}{RGB}{255, 184, 28}

\begin{document}

% =======================================================
% INSUMO 1: GUÍA DE CLASE PARA EL DOCENTE
% =======================================================
\begin{center}
    \Large\textbf{\textcolor{ucbblue}{GUÍA DE CLASE PARA EL DOCENTE (LESSON PLAN)}}\\
    \vspace{0.2cm}
    \normalsize\textbf{Código:} IMT-342-GP-0201-V2
\end{center}
\vspace{0.3cm}

\noindent\textbf{Unidad I:} Morfología, Geometría del Espacio y Cinemática Directa \\
\textbf{Tema:} Grupo $SO(3)$, Dualidad Activa/Pasiva y Composición de Rotaciones \\
\textbf{Duración:} 120 minutos reales \\
\textbf{Alineamiento:} CDIO (Diseñar/Implementar) + ABET SO-1 (Análisis de Ciencias e Ingeniería)

\vspace{0.5cm}
\noindent\textbf{\large Estructura de la Sesión y Gestión del Tiempo (120 min)}
\vspace{0.2cm}

\noindent
\begin{tabularx}{\textwidth}{@{} l >{\raggedright\arraybackslash}p{3cm} X X @{}}
\toprule
\textbf{Tiempo} & \textbf{Fase} & \textbf{Actividad Docente y Estudiantil} & \textbf{Justificación Pedagógica} \\ \midrule
16:00 - 16:15 & Activación y \newline Problema Físico & Demostración con objeto físico: Rotar un vector (Activo) vs. Mover la cámara (Pasivo). & \textbf{Anclaje Físico:} Elimina ambigüedad entre mover un eslabón o cambiar de marco. \\ \addlinespace
16:15 - 17:00 & Cátedra \newline Analítica & Deducción de $SO(3)$. Demostración de $p' = R \cdot p$ vs $p_A = {}^{A}R_B \cdot p_B$. Deducción $R^T R = I$. & \textbf{Rigor Matemático:} Fundamentación analítica, evitando el enfoque de prueba y error. \\ \addlinespace
17:00 - 17:35 & Composición de \newline Rotaciones & Análisis de Ejes Fijos (Pre-multiplicación) vs. Ejes Móviles (Post-multiplicación). & \textbf{Alineamiento PFI:} Previene errores críticos al multiplicar matrices D-H. \\ \addlinespace
17:35 - 17:55 & Taller en Python \newline (NumPy) & Código en Ubuntu: Verifican $R^T R = I$, conservación de norma y transformaciones. & \textbf{Implementación:} Principios 2 y 4 de la metodología CDIO en tiempo real. \\ \addlinespace
17:55 - 18:00 & Cierre y \newline Asignación & Entrega de Guía Práctica 1 y asignación de lectura sobre Euler para el jueves. & \textbf{Progresión Lógica:} Prepara el terreno minimizando la sobrecarga cognitiva. \\ \bottomrule
\end{tabularx}

\vspace{0.5cm}
\begin{tcolorbox}[colback=ucbblue!5!white, colframe=ucbblue, title=Ajuste de Diseño Didáctico]
Se ha decidido \textbf{diferir} la parametrización de Euler y Roll-Pitch-Yaw a la Sesión 2. Esta clase se concentra de forma exclusiva en la dualidad de la matriz en $SO(3)$ (operador vs proyector) para asentar las bases algebraicas antes de abordar el \textit{Gimbal Lock}.
\end{tcolorbox}

\newpage
% =======================================================
% INSUMO 3: GUÍA DE TRABAJO PRÁCTICO
% =======================================================
\begin{center}
    \Large\textbf{\textcolor{ucbblue}{UNIVERSIDAD CAT\'OLICA BOLIVIANA ``SAN PABLO''}}\\
    \large\textbf{SEDE TARIJA}\\
    \vspace{0.2cm}
    \normalsize Departamento de Ciencias Exactas e Ingenierías | Carrera de Ingeniería Mecatrónica\\
    Asignatura: Robótica (IMT-342) | Gestión: II-2026
\end{center}

\vspace{0.4cm}
\begin{center}
    \Large\textbf{Hoja de Trabajo Práctico Nº 1}\\
    \large\textit{Geometría de Rotación, Dualidad de $SO(3)$ y Composición}
\end{center}
\vspace{0.4cm}

\section{Principios Obligatorios de la Práctica (CDIO)}
\begin{itemize}[label=\textcolor{ucbblue}{$\blacktriangleright$}, leftmargin=*]
    \item \textbf{Fundamentación Teórica:} Demostración matemática manuscrita paso a paso.
    \item \textbf{Implementación:} Script modular en Python utilizando \texttt{NumPy}.
    \item \textbf{Validación:} Verificación numérica de la isometría ($\left\|p'\right\| = \left\|p\right\|$) y ortogonalidad ($R^T R = I$).
    \item \textbf{Evidencia:} Jupyter Notebook entregado en el repositorio Git del equipo.
\end{itemize}

\section{Ejercicios de Desarrollo Obligatorio}

\subsection*{Ejercicio 1: Análisis de la Dualidad Activa vs. Pasiva (Teoría)}
Dado el vector $p = [1, 2, 0]^T$ y la matriz de rotación canónica $R_z(90^\circ)$:
\begin{itemize}
    \item \textbf{Caso A (Transformación Activa):} Calcule $p' = R_z(90^\circ) \cdot p$. Grafique en papel milimetrado el vector original $p$ y el vector resultante $p'$ dentro del \textit{mismo} sistema de referencia $\{A\}$.
    \item \textbf{Caso B (Transformación Pasiva):} Suponga que $p$ representa las coordenadas de un punto respecto a un marco móvil $\{B\}$ ($p_B = [1, 2, 0]^T$). Si el marco $\{B\}$ está rotado $90^\circ$ respecto al fijo $\{A\}$ alrededor de $Z_A$, calcule $p_A = {}^{A}R_B \cdot p_B$. Grafique los dos marcos de referencia $\{A\}$ y $\{B\}$ y demuestre geométricamente por qué $p_A$ describe el mismo punto físico.
\end{itemize}

\subsection*{Ejercicio 2: Composición de Rotaciones (Ejes Fijos vs. Móviles)}
Un marco de referencia móvil $\{B\}$, inicialmente alineado con el marco fijo $\{A\}$, realiza la siguiente secuencia de movimientos:
\begin{enumerate}
    \item Gira $45^\circ$ alrededor del eje $X_A$ \textbf{fijo}.
    \item Gira $90^\circ$ alrededor del eje $Z_B$ \textbf{móvil}.
    \item Gira $-45^\circ$ alrededor del eje $Y_A$ \textbf{fijo}.
\end{enumerate}
\textbf{Tarea:} Obtenga analíticamente la matriz de rotación final $^{A}R_B$. Escriba la ecuación de multiplicación indicando explícitamente cuáles términos se pre-multiplican y cuáles se post-multiplican.

\subsection*{Ejercicio 3: Implementación y Validación en Python (NumPy)}
Escriba un Jupyter Notebook en Ubuntu que:
\begin{enumerate}
    \item Defina las funciones \texttt{rot\_x(deg)}, \texttt{rot\_y(deg)} y \texttt{rot\_z(deg)}.
    \item Calcule numéricamente la matriz $^{A}R_B$ del Ejercicio 2.
    \item Compruebe mediante la norma de Frobenius que $\left\| R^T R - I \right\|_F < 10^{-15}$.
    \item Tome un vector $p_B = [2, -1, 4]^T$, compute $p_A = {}^{A}R_B \cdot p_B$ y demuestre por código que $\left\| p_A \right\| - \left\| p_B \right\| = 0$.
\end{enumerate}

\section{Criterios de Evaluación y Entrega}
\begin{itemize}
    \item \textbf{Entregable:} Jupyter Notebook en la rama \texttt{feature/tp1-rotation} del repositorio.
    \item \textbf{Fecha Límite:} Jueves 13 de Agosto, 16:00 h.
\end{itemize}

\vspace{0.2cm}
\noindent\textbf{Rúbrica de Evaluación (Alineamiento ABET/CDIO):}
\begin{table}[h!]
    \centering
    \begin{tabular}{@{} l c @{}}
    \toprule
    \textbf{Criterio de Evaluación} & \textbf{Ponderación} \\ \midrule
    Demostración Teórica Manuscrita \textbf{(ABET SO-1)} & 40\% \\
    Código en Python y Modularidad \textbf{(CDIO)} & 40\% \\
    Validación del Error Numérico e Isometría \textbf{(ABET SO-6)} & 20\% \\ \bottomrule
    \end{tabular}
\end{table}

\end{document}